\chapter{Reactions Driven by Electron Transfer}

\begin{kaichang}
Find a rusty iron nail: on a balcony, an old bicycle, or a screw in outdoor exercise equipment. Compare it with a new nail, looking at them and weighing them in your hands. The old one wears a loose reddish-brown crust; the new one is smooth, gray, and metallic.

Now guess three things. First, does iron become heavier or lighter when it rusts? Many people instinctively say, ``It has rotted away, so of course it is lighter.'' Second, did this reddish-brown crust form from iron ``going bad,'' or was it ``stuck on'' from outside? Third, a cut apple turns yellow-brown after a while, as does a cut potato. Is there a connection between what happens to metal and what happens to fruit?

Write down your guesses. By this chapter's end, you will discover that rusting, apple browning, burning matches, bleaching colors, and your breathing right now all enact the same microscopic drama: \textbf{electrons pass from one atom's hands into another's}.
\end{kaichang}

\section{Visible Transformations}

Start with observation, without rushing to explain. Set these events side by side:

\begin{itemize}[itemsep=2pt]
\item Iron nails develop reddish-brown rust spots within days in damp air, but rust slowly when kept dry or coated in oil.
\item A match strikes a phosphorus-coated surface and flares brightly with a scratch. After burning, it leaves a short black charcoal stick and gives off white smoke.
\item An apple's cut surface turns from white to yellow-brown after half an hour exposed to air, but changes much more slowly in saltwater or lemon juice.
\item A white cloth stained with blue-black ink visibly loses its color when dilute bleach is applied.
\item In a laboratory demonstration, bright copper wire placed in silver nitrate solution slowly ``grows'' silvery fuzz, while the initially colorless solution turns blue.
\end{itemize}

Five events, five color changes, seemingly unrelated. Yet they share one suspect: \textbf{oxygen or electrons}. Rust contains oxygen, burning requires oxygen gas, and apple browning needs air. What about copper ``growing silver''? No oxygen gas participates in that solution reaction. And if you weigh everything after a match burns, its smoke and ash together are \textbf{heavier} than the original matchstick. Where did the extra mass come from?

Chemists in the 1770s were stuck on this question. What trapped them was an old, plausible-sounding theory, which we will meet in the evidence story. First, zoom down to the particles yourself.

\section{Zoom In: Electrons Change Owners}

Recall Chapter 9's picture: electrons live outside atomic nuclei, and the outermost ones determine an element's temperament. Chapter 12 explained that metals such as sodium have outer electrons that ``do not sit securely'' and are easily lost; Chapter 13 described halogens such as chlorine as ``grabbing any electron in sight.'' Chapter 18 showed what follows the handover: positively charged sodium ions and negatively charged chloride ions embrace through electrostatic attraction to form salt.

Now apply this picture to burning magnesium ribbon. Magnesium burns vigorously in air with a dazzling white light---used in fireworks and flash lamps---forming white magnesium oxide powder. Two things happen microscopically: every magnesium atom gives up two electrons, and every oxygen atom accepts two; then the charged particles join. Neither magnesium nor oxygen has ``gone bad'' into a different element: their proton counts are unchanged. Only \textbf{the electrons have changed owners}.

Consider copper wire in silver nitrate again. A silver ion, \chem{Ag^+}, meeting a copper atom ``collects'' one electron, becomes a neutral silver atom, and deposits on the wire: that is the silvery fuzz. Copper losing electrons becomes \chem{Cu^{2+}} and leaves the surface for the solution, coloring it blue. Notice that \textbf{no oxygen participates}, yet this belongs to the same reaction class as burning magnesium: both transfer electrons.

Chemists therefore gave these reactions one name: \textbf{oxidation--reduction reactions}, or redox reactions. There is just one criterion:

\begin{itemize}[itemsep=2pt]
\item A particle that \textbf{loses} electrons is \textbf{oxidized};
\item A particle that \textbf{gains} electrons is \textbf{reduced}.
\end{itemize}

The name is rather unfortunate. ``Oxidation'' comes from oxygen, historically the first recognized ``expert electron grabber,'' whose name was applied to the whole class. But names are historical fossils. Silver ions oxidize copper without oxygen; hydrogen burns quietly in chlorine to form hydrogen chloride, \chem{HCl}, without oxygen, and that too is redox: hydrogen loses electrons and chlorine gains them. \textbf{Redox does not mean gaining or losing oxygen}: this is our first firm boundary. Track electrons, not oxygen's presence or absence.

There is another accounting rule: electrons neither vanish nor drift homeless through a solution. In a reaction, the total lost in one place must equal the total gained elsewhere. Oxidation and reduction always \textbf{occur together}, like two sides of a coin. You cannot stop after saying one substance ``was oxidized''; you must also ask what was reduced.

This introduces two familiar terms: the one that \textbf{gives away} electrons and reduces another substance is the \textbf{reducing agent}; the one that \textbf{takes in} electrons and oxidizes another is the \textbf{oxidizing agent}. When magnesium burns, magnesium is the reducing agent, itself oxidized, and oxygen is the oxidizing agent, itself reduced. Confusing? Remember this image: an \textbf{agent} is the actor, named for what it does to \textbf{someone else}, while the opposite happens to itself. Oxygen routinely oxidizes other substances, so it is an oxidizing agent.

Identify both roles in the match you struck. Its head contains the oxidizer potassium chlorate, \chem{KClO_3}, an ``electron buyer'' stockpiling oxygen, along with reducing sulfur and phosphorus compounds. The red phosphorus coating on the box provides another reducing agent. Striking produces frictional heat, supplying Chapter 29's activation energy to cross the energy mountain. Electrons then relocate between oxidizing and reducing agents, releasing enough heat to sustain the flame. A tiny match head is a packaged redox reaction awaiting your trigger. Safety matches separate the oxidizer from some of the reducing agent, placing them on the head and the box's striking surface. This uses the rule ``both partners must meet'' to prevent fires. Redox knowledge is hidden in the small question of why a match must be struck to light.

\section{Electron Accounting: Oxidation Numbers}

Electron transfer sounds clear, but complex molecules complicate it. In potassium permanganate, \chem{KMnO_4}, or sulfuric acid, \chem{H_2SO_4}, electrons around manganese or sulfur are shared in covalent bonds (Chapter 19), rather than whole electrons truly moving house. How do we keep those accounts?

Chemists introduce an accounting tool: the \textbf{oxidation number}, also called oxidation state. Its central rule is a ``pretend settlement'': award every shared electron in each covalent bond to the more electronegative atom, the stronger electron attractor of Chapter 11. Then count how many more or fewer electrons each atom ``owns'' than when neutral. Extra electrons give a negative number, missing electrons a positive one. Common rules are:

\begin{itemize}[itemsep=2pt]
\item Atoms in elemental substances have oxidation number 0: \chem{Fe}, \chem{O_2}, and \chem{Cu} all count as 0.
\item A monatomic ion's oxidation number equals its charge: \chem{Na^+} is $+1$ and \chem{Cl^-} is $-1$.
\item In compounds, oxygen usually counts as $-2$, except in peroxides and other exceptions; hydrogen usually counts as $+1$, except in metal hydrides.
\item All atomic oxidation numbers sum to 0 in a neutral molecule, or to the ion's charge in an ion.
\end{itemize}

Calculate potassium permanganate: potassium is $+1$, four oxygens at $-2$ give $-8$, and the total must be 0, so manganese must be $+7$. Next, consider burning glucose, respiration's main character. In \chem{C_6H_{12}O_6}, hydrogen is $+1$ and oxygen $-2$, so the six carbons must total $0-(+12)-(-12)=0$, averaging 0 each. After burning to \chem{CO_2}, carbon is $+4$. Every carbon has ``lost'' four electrons: glucose has been oxidized. This is how your body obtains energy from sugar, as Chapters 52 and 53 explain.

Here is the second firm boundary: \textbf{oxidation numbers are accounting tools, not little labels attached to atoms that can be measured directly}. Manganese in \chem{MnO_4^-} does not really carry charge $+7$, which would mean seven electrons had moved away completely. Manganese and oxygen actually share electrons; our bookkeeping simply awards more to oxygen. Real charge distributions require quantum calculations and precise experiments such as X-ray spectroscopy to approximate, often yielding untidy fractions. Oxidation numbers matter not because they are ``real,'' but because they are \textbf{useful}: who rises, who falls, and by how much become immediately clear, especially when balancing equations. Think of them as account codes in double-entry bookkeeping, not cash in a safe.

We can now restate our definitions without relying on whether electrons literally move house: increasing oxidation number $=$ oxidation; decreasing oxidation number $=$ reduction. Carbon rises from 0 to $+4$, while manganese falls from $+7$ to $+2$ in a permanganate decolorization experiment. That is how the ledger works.

\begin{qiaoliang}
How many electrons move when an iron nail rusts? Take a nail of about \ud{5}{g}. Iron's molar mass is approximately \ud{56}{g/mol}, so it contains $5/56 \approx 0.09\ \mathrm{mol}$ of iron atoms. In rust's main component, iron is $+3$, so each atom gives up about three electrons. Completely rusting the nail therefore releases about $0.27\ \mathrm{mol}$ of electrons. Multiply by \ud{6.02\times 10^{23}}{mol^{-1}} to obtain about $1.6\times 10^{23}$ electrons---sixteen quintillion electrons passing one by one from iron atoms to oxygen atoms. Electrons moved each hour in your kitchen from gas, or through air and food, number just as astronomically. Chemical changes sound gentle, but their electron accounts are always torrents.
\end{qiaoliang}

\section{Split It in Two: Half-Reactions}

The best way to balance the books is to \textbf{separate oxidation and reduction into half-reactions}. Return to copper and silver nitrate. On the copper side:

\[\chem{Cu} \rightarrow \chem{Cu^{2+}} + \chem{2e^-}\]

On the silver side:

\[\chem{Ag^+} + \chem{e^-} \rightarrow \chem{Ag}\]

Every two electrons released by copper need two silver ions to accept them. Multiply the entire second half-reaction by 2 and add; the electrons cancel exactly:

\[\chem{Cu} + \chem{2Ag^+} \rightarrow \chem{Cu^{2+}} + \chem{2Ag}\]

This is the \textbf{half-reaction method}. Its one rule is \textbf{electron conservation}: electrons handed over on one side must equal those accepted on the other. Add Chapter 26's atom conservation and check equal total charges on both sides, and the redox equation is balanced. You know exactly where every electron came from and went, without guessing coefficients.

The method has a hidden benefit worth planting here: if half-reactions can be written separately, can they happen in separate places, with copper losing electrons here and silver ions accepting them there, linked by an \textbf{external circuit}? Yes. Then it is no longer merely a beaker reaction: it is a battery. Chapter 34 devotes itself to this.

\section{A Debate Ended by a Balance}

\begin{zhengju}
The eighteenth century's popular combustion theory was ``phlogiston'': combustible substances contained phlogiston, and burning released it into the air. It explained flames and smoke and seemed plausible. But one problem became increasingly glaring: the ``calx'' left by burned metal, now called an oxide, was \textbf{heavier than the original metal}. How could losing something increase weight? Supporters patched the theory: phlogiston had ``negative weight.'' Once a theory relies on negative weight for survival, trouble is near.

In the 1770s, the French chemist Lavoisier brought the balance center stage. He heated tin and lead in sealed containers and weighed the total: unchanged. Open the container, let air rush in, and weigh again: the added mass exactly matched the metal's gain on becoming calx. He also decomposed mercury calx by heating and collected a gas whose mass balanced the account, and in which candles burned exceptionally vigorously. He named it oxygen. The conclusion was clear: burning did not release phlogiston, but \textbf{combined a substance with oxygen}; the extra mass came from air. Total mass before and after reaction in a sealed container stayed constant, one of the classic experimental foundations for Chapter 26's conservation of mass.

It is worth noting that Lavoisier was not the first to prepare the gas. Britain's Priestley and Sweden's Scheele obtained oxygen earlier. Priestley even found that candles burned brightly and mice thrived in it, but interpreted it through phlogiston theory as ``dephlogisticated air.'' Rather than helping combustion, it supposedly ``absorbed phlogiston particularly well.'' The same observations told wholly different stories in different theoretical frameworks. This contrast offers one of science history's best lessons: \textbf{experimental data do not come with their own explanations}. What you see matters, and so do the concepts used to organize it.

This tells only half the story. Lavoisier answered where mass came from, but not what combining truly meant. He could not: the electron would not be discovered for over another century, in 1897. Once twentieth-century chemists developed the electron-transfer picture, they saw that oxygen starred in combustion simply because it was good at taking electrons. Burning, rusting, and respiration all shared electron transfer. ``Oxidation'' was therefore \textbf{redefined} from ``combining with oxygen'' to ``losing electrons.'' Scientific concepts often work this way: the word remains while its meaning changes. Learn a scientific term by what it means now, not what its wording seems to suggest.
\end{zhengju}

\begin{shiyan}
\textbf{Why Iron Nails Rust: Three Comparisons}

Materials: three clean, rust-free iron nails, sanded bright; three clear cups; tap water; salt; cooking oil; desiccant, such as the small packets in snack bags; plastic wrap.

Procedure: (1) Half-fill Cup 1 with tap water and add a nail, leaving it open with the nail half submerged and half exposed to air. (2) Put the same amount of water in Cup 2, dissolve a spoonful of salt, and add a nail. (3) Put no water in Cup 3; add a nail and two desiccant packets, and seal tightly with plastic wrap. If cooking oil is available, add a fourth cup: cover the water with an oil layer before adding the nail, to explore ``water without air.''

What to observe: inspect once a day for three days. Record which cup first develops reddish-brown spots and where rust appears: near the waterline or at the bottom?

Expected results and questions: Cup 2 usually rusts fastest; Cup 3 hardly rusts. This shows that iron rusting \textbf{requires both water and oxygen}, and salt greatly accelerates it: electrons ``travel'' more easily in saltwater. Why is rust heaviest near Cup 1's waterline? Hint: both oxygen and water are most abundant there.

Safety: the experiment itself is safe, but used nail edges may be sharp; hold them through paper. Promptly clean up sanding debris and keep it out of your eyes.
\end{shiyan}

\section{Where This Ledger Can Mislead}

Oxidation numbers are useful, but remember their ``pretend settlement'' basis. Three situations expose it.

First, they often do not give real charges. Carbon in \chem{CO_2} counts as $+4$ but does not carry charge $+4$: the real charge displacement is much smaller and continuously distributed over atoms and bonds. Using oxidation numbers to calculate actual electrostatic forces gives absurd results.

Second, average oxidation numbers can be fractional. Iron's average in magnetite, \chem{Fe_3O_4}, is $+8/3$. There is no ``eight-thirds-valent'' atom. The crystal contains two types of iron, some $+2$ and some $+3$, averaging to a fraction. The ledger gives an average, but the shelves hold two products.

Third, determining whether a reaction is redox requires checking each element, not going by appearances. Calcium carbonate decomposing into calcium oxide and carbon dioxide looks dramatic, but every element's oxidation number stays unchanged: no electron transfer, no redox. Likewise, Chapter 32's acid--base neutralization is not redox: protons transfer, electrons do not. Dividing chemical reactions into those driven by electron transfer and those without it is far more reliable than asking whether oxygen or flames are involved.

\begin{wuqu}
``Oxidation means reacting with oxygen gas.'' This is the unfortunate name's greatest misconception. Counterexamples abound: hydrogen burns in chlorine to make hydrogen chloride, with no oxygen, yet hydrogen is oxidized from 0 to $+1$. Copper in silver nitrate becomes \chem{Cu^{2+}} with no oxygen involved. Iron ions in your body can switch between $+2$ and $+3$ without direct participation by oxygen gas. The sole criterion is electrons or oxidation number: losing electrons and rising oxidation number mean oxidation. Oxygen is merely the most famous ``electron buyer.'' There is an opposite misconception too: reduction does not mean ``returning to the original substance,'' but \textbf{gaining electrons}. When iron ore is reduced to iron, iron ions gain electrons, and oxygen is what gets removed.
\end{wuqu}

\section{Back to the Iron Nail}

Now answer the opening questions.

First, rust is \textbf{heavier} than the iron it came from: the extra is oxygen combined from air. The nail may seem lighter only because its loose rust crust flakes off. Collect and weigh all the flakes, and the account balances, showing a gain. Lavoisier's balance established this over two hundred years ago.

Second, the reddish-brown crust is neither iron ``going bad'' nor dust stuck on from outside. Iron atoms lose electrons, become iron ions, and combine with oxygen and water to form hydrated iron oxide. The iron element remains, but its electron state and neighbors change: the same identity card, different roles, as Chapter 10 put it.

Third, apples and nails perform the same drama. Phenolic substances in a cut apple meet atmospheric oxygen and, helped by enzymes released from cells (Chapter 49), oxidize into yellow-brown products. You see slow, enzyme-accelerated oxidation. Saltwater and lemon juice work by interfering with enzymes or changing acidity, slowing oxidation. Rusting is oxidation too, but without enzyme help, proceeding at its own modest pace. Burning is oxidation fast enough to glow; rusting and browning are almost imperceptibly slow oxidation; your respiration, coming next, is oxidation whose speed life \textbf{controls precisely}.

\section{From Rust to Respiration: Transfer Everywhere}

Zoom out, and electron transfer appears as one of this planet's busiest businesses.

Smelters essentially use strong reducing agents to ``push electrons back'' into metal ions in ores. In a blast furnace, carbon monoxide gives electrons to iron ions in iron oxide, reducing them to molten iron that flows out: this is how Chapter 15's iron is ``redeemed'' from ore. Bleaches and disinfectants work in the opposite direction. Sodium hypochlorite is a fierce oxidizer that forcibly oxidizes and disrupts pigment molecules and vital microbial molecules, removing colors and killing bacteria. This is why bleach must never be mixed with acidic toilet cleaner: it releases chlorine gas, one of redox chemistry's greatest household dangers.

Keep the accounts for daily kitchen combustion too. Methane, \chem{CH_4}, the main component of natural gas, burns:
\[\chem{CH_4} + \chem{2O_2} \rightarrow \chem{CO_2} + \chem{2H_2O}\]
Check the oxidation numbers. Methane's carbon is $-4$: hydrogen is $+1$, and the more electronegative carbon is awarded shared electrons. Carbon dioxide's carbon is $+4$, so one carbon atom gives up eight electrons. Oxygen falls from 0 to $-2$, and two oxygen molecules accept eight electrons altogether, balancing the books. The blue flame is light from this electron marketplace's countless transactions each second. Stove ventilation and carbon monoxide from incomplete combustion are safety lessons grounded in the same accounts: without enough oxidizer, carbon can hand over only ``half'' its electrons, stopping at a toxic intermediate.

The most important redox reaction is happening in your body right now. Carbon in the glucose you eat has average oxidation number 0; carbon in exhaled carbon dioxide is $+4$. From mouth to lungs, every carbon gives up four electrons: \textbf{exactly the same overall accounting} as burning glucose. The difference is that a bonfire lets electrons ``rush in together,'' wasting energy as light and heat, while cells build an ``electron transport chain'' dozens of steps long. Electrons proceed in small steps, and each small release of energy is collected and stored in ATP. Respiration is more than ``inhale oxygen, exhale carbon dioxide.'' Oxygen stands at the chain's endpoint, finally accepting electrons and forming water, as Chapters 52 and 53 will unpack. Photosynthesis reverses the overall ledger, using sunlight to force electrons back into carbon dioxide and remake sugar. Its oxygen comes from split water molecules, not carbon dioxide: Chapter 54 returns to this firm rule for the whole book, so remember the conclusion now.

From nails to blast furnaces, apples to lungs, redox is a thread through both inorganic and living worlds. It also has an enticing branch already mentioned: half-reactions separable on paper can be separated in space. Put oxidation and reduction in different ``rooms,'' force electrons along a chosen external circuit, and they can work along the way: light bulbs, drive motors, charge phones. Chemical change becomes electric current. Next chapter, we will build a battery.

\begin{tiaozhan}
Try balancing a classic reaction with half-reactions: purple-red potassium permanganate solution loses color on meeting oxalic acid, \chem{H_2C_2O_4}, found in spinach and tea. Start with manganese: \chem{MnO_4^-} becomes \chem{Mn^{2+}}, falling from $+7$ to $+2$ and accepting five electrons. Balance the excess oxygen on the left with \chem{H^+} and water:
\[\chem{MnO_4^-} + \chem{8H^+} + \chem{5e^-} \rightarrow \chem{Mn^{2+}} + \chem{4H_2O}\]
For carbon, oxalic acid's $+3$ becomes $+4$ in \chem{CO_2}, so the two carbons give up two electrons:
\[\chem{H_2C_2O_4} \rightarrow \chem{2CO_2} + \chem{2H^+} + \chem{2e^-}\]
The least common multiple is 10: multiply the first equation by 2 and the second by 5, add, cancel electrons, and simplify to obtain the balanced result. Count every atom and total charge on both sides: the books should balance. Another interesting phenomenon is \textbf{disproportionation}, in which some atoms of the same element in a molecule rise in oxidation number while others fall. Hydrogen peroxide decomposition is an example: oxygen goes from $-1$ to $-2$ and 0. It serves as both its own oxidizing and reducing agent. Skipping this section does not disrupt the main thread.
\end{tiaozhan}

\section*{Try It Yourself}

\begin{sikaoti}
\item Check the accounts: which are redox processes? (1) Water freezing; (2) burning natural gas, methane \chem{CH_4}; (3) neutralizing hydrochloric acid with sodium hydroxide; (4) photosynthesis in green plants. For each redox process, identify which element is oxidized and which reduced.
\item Draw a particle diagram of copper wire in silver nitrate solution. Label copper atoms, copper ions, silver ions, and silver atoms, show electron routes with arrows, and write both half-reactions beside it. Note that ball sizes and colors are human-made representations.
\item Use oxidation numbers: what is sulfur's oxidation number in \chem{SO_2}, \chem{SO_4^{2-}}, and \chem{H_2S}? When sulfur near a volcanic vent is oxidized from \chem{H_2S} to \chem{SO_2}, how many electrons does each sulfur atom give up?
\item A household has both ``84 disinfectant,'' sodium hypochlorite, an oxidizer, and toilet cleaner containing hydrochloric acid. Look up their ingredients and explain why mixing releases chlorine. From a redox perspective, what happens to chlorine? Research only: \textbf{never} mix them to test it.
\item Draw an electron-flow diagram with two routes for glucose burning in a bonfire and glucose breaking down in cells through respiration, sharing endpoints but differing in path. Mark which releases energy all at once and which splits release into small steps. Why can cells ``save money'' with the same overall accounts?
\item Someone says, ``Rusting oxidizes iron, so rust is iron gone bad; just keep oxygen out and you never need to worry again.'' Evaluate this using the chapter's model. Is the first part correct? Is the second idea feasible in engineering---consider ships, bridges, and food packaging? Are there cases where ``letting iron rust is fine''?
\end{sikaoti}
